\section{Introduction}
\label{sec:intro}

Pretrained language models encode rich semantic knowledge in their embedding layers and residual streams. A natural question arises: if we systematically remap \emph{which} token IDs correspond to \emph{which} surface forms---applying a bijective permutation $\perm: \vocab \to \vocab$ to the input---can the model's internal representations adapt? After processing enough permuted text, does the residual stream converge to something resembling its normal state, ``modulo the shuffle''?

This question matters for several reasons. First, it probes how deeply a transformer's ``understanding'' is tied to specific token identities versus distributional structure. Second, it has practical implications for adversarial robustness: if a simple bijective shuffle can destroy a model's predictions even though the statistical structure of the input is preserved, this reveals a fundamental rigidity in how transformers process language. Third, it informs the growing literature on cipher-based jailbreaks~\citep{yuan2023gpt4cipher}, where attackers encode harmful prompts via substitution ciphers to bypass safety filters.

{\bf What is missing in existing work?} Recent work by \citet{fang2025iclciphers} demonstrates that models \emph{can} learn bijective ciphers from in-context demonstrations, with the logit lens showing progressive cipher decoding at deeper layers. However, their analysis focuses on the \emph{with-demonstrations} setting and measures only top-token predictions, not the full geometry of the residual stream. No prior work has systematically measured how the residual stream's \emph{structure}---its direction, magnitude, and representational content---changes under bijective token permutation in the zero-shot setting, or whether there exists a linear transformation that maps ciphered hidden states back to clean ones.

We address this gap directly. We extract residual stream activations at every layer of \gpttwo (12 layers) and \pythia (24 layers) under three conditions: clean text, bijectively ciphered text, and random token replacement. We compare these using cosine similarity, a permutation-adjustment matrix, the logit lens, and centered kernel alignment (CKA). Our key innovation is the \emph{permutation-adjusted similarity} test: if the residual stream is truly ``normal modulo the shuffle,'' then transforming ciphered hidden states by the cipher's inverse in embedding space should dramatically increase similarity to clean states.

{\bf Our results are clear}: the model does not recover. The permutation adjustment makes similarity \emph{worse}, not better. The high cosine similarity observed at later layers in \gpttwo (0.873 for full shuffle) is an architectural artifact---\layernorm drives \emph{all} inputs, including random tokens (0.858), toward similar representations. The logit lens confirms the model never predicts cipher-consistent next tokens at any layer. In \pythia, which exhibits weaker architectural convergence, ciphered representations sit between clean and random, suggesting partial distributional extraction without genuine cipher decoding.

We make the following contributions:
\begin{itemize}[leftmargin=*,itemsep=0pt,topsep=0pt]
    \item We conduct the first systematic analysis of residual stream geometry under bijective token permutation, testing whether pretrained transformers can internally remap a cipher without demonstrations.
    \item We introduce a permutation-adjustment test that directly evaluates the ``normal modulo the shuffle'' hypothesis, finding it decisively refuted.
    \item We identify \layernorm-driven architectural convergence as a confound that inflates cosine similarity between \emph{any} two inputs, cautioning against using raw cosine similarity as a measure of representational understanding.
    \item We validate our findings across two model architectures (\gpttwo and \pythia), four shuffle rates, and five complementary analysis methods.
\end{itemize}
